%%%%%%%%%%%%%%%%%%%%%%%%%%%%%%%%
\chapter{はじめに}\label{Introduction}
%%%%%%%%%%%%%%%%%%%%%%%%%%%%%%%%

\section{問題と背景}
%%%%%%%%%%%%%%%%% [[[[問題と背景]]]] %%%%%%%%%%%%%%%%%%%%%%%%%%%%%%%%
自然界における様々な現象を数学的に記述し，その問題を計算機を用いて数値的に解くことを数値計算と呼ぶ．
数値計算で用いる計算機は，実数をIEEE754 Standard で定められた浮動小数点数に近似して演算を行う．
この手法は非常に高速な計算が可能であるが，一方で得られる結果は丸め誤差を含む近似値となる．
また，数値計算では丸め誤差だけではなく，離散化誤差や打切り誤差などが生じる．
例として，Newton法は理論上，無限回の演算を行うことで真の解を求める手法であるが，
計算機では無限回の計算は不可能であるため，有限回の計算で打ち切る．結果として誤差が生じ，得られる解は近似解となる．
このような誤差を伴う計算結果に対して，その信頼性を確認する手段として精度保証付き数値計算が存在する．
精度保証付き数値計算とは，対象とする数学的問題の真の解\(x^*\) と近似解\(\hat{x}\)に対し，
計算機を用いて数学的に厳密な誤差の上界を定量的に得る手法である．
非線形偏微分方程式にも精度保証付き数値計算を用いることができるが，その方程式が無限次元問題の場合，
計算機上でそのまま計算をすることができない．
そのため，スペクトル法や有限要素法を用いて離散化し，有限次元の問題に近似したうえで計算機を用いて計算する．
この時，離散化誤差への考慮が必要となる．
また，解析対象となる問題がそもそも真の解をもつかどうか分からない場合も少なくない．
中尾\cite{artnakao}により，関数解析やSobolev空間論の理論を精度保証付き数値計算に導入し，
無限次元問題の解の存在性や一意性の範囲を定量的に特定する研究がなされた．
以後，Plum\cite{artplum}，大石\cite{artoisi}をはじめとする多くの研究者により研究が進められている．
本論文では，フィッツヒュー・南雲方程式
\begin{equation}\label{eq:問題}
  \left\{ \,
      \begin{aligned}
          -{\epsilon}^2\Delta u &= f(u)-\delta v \quad &\text{in}\;\;\Omega\\
          -\Delta v &=u-rv \quad &\text{in} \;\;\Omega\\
          u &= v=0 \quad &\text{on}\;\;\partial\Omega\\
      \end{aligned}
  \right.
\end{equation}
のDirichlet境界値問題の解に対して\(\tilde{\mathcal{F}}\)ノルムを用いた精度保証法を確立することが目的である．
本論文では(\ref{eq:問題})の非自明解を精度保証付きで求めることを目標とする．

\section{本論文の構成}
%%%%%%%%%%%%%%%%% [[[[本論文の構成]]]] %%%%%%%%%%%%%%%%%%%%%%%%%%%%%%%%
本論文は以下のように構成されている．
２章ではBanach空間，Hilbert空間，作用素などの必要な知識についての説明をしている．
３章では重み付きの内積を定義した空間でNewton-Kantorovichの定理\cite{bookoisi}を適用し，
逆作用素ノルムを無限次元固有値問題を通じて評価する従来の手法をフィッツヒュー・南雲方程式に適用する手法について述べている．
4章では新しくフィッツヒュー・南雲方程式に特化した内積を定義した空間を設定した手法を提案し，
フィッツヒュー・南雲方程式ののDirichlet境界値問題の解に対する．精度保証付き数値計算の手法を述べている．
５章では実際に数値実験を行った結果を示し，６章で本論文のまとめと今後の課題を示す．

